$E_{\text{core}}$ distributions in real data as well as in MC were investigated. MC shows that $\pi^\pm$ and $K^\pm$ mostly deposit $E_{\text{core}} > 0.25$ in PbSc and $E_{\text{core}} > 0.35$ in PbGl, while real data distributions have many $E_{\text{core}}$ deposits in lower regions. 
This, apart from $p+\bar{p}$ deposits, indicates significant contamination in these regions. 
In previous analyses (\href{https://phenix-intra.sdcc.bnl.gov/phenix/WWW/p/info/an/375/}{AN375} and \href{https://phenix-intra.sdcc.bnl.gov/phenix/WWW/p/draft/berd/Sergei/AN_AuAu_KStar.pdf}{AN1504}) $E_{\text{core}}$ cut was shown to significantly decrease background contamination in $m^2$ distribution for charged hadron identification in EMCal which shows its effectiveness.
Therefore $E_{\text{core}} > 0.25$ in PbSc and $E_{\text{core}} > 0.35$ in PbGl were implemented as cuts in real data and MC.
